% ---------- 6. Big philosophical questions ----------
\begin{frame}{Twenty-four questions about minds and machines}
\begin{sfig}[0.7]
\begin{tikzpicture}
  \node[card=tdBlue,text width=5.0cm,minimum height=3.1cm,anchor=north west] at (0,0) {\textbf{Mind and consciousness}\\[2pt]\shortstack{\scriptsize 1 hard problem\\ 2 Mary's room\\ 3 zombies\\ 4 Nagel's bat\\ 5 functionalism vs biological naturalism\\ 6 intentional stance}};
  \node[card=tdGreen,text width=5.0cm,minimum height=3.1cm,anchor=north west] at (5.4,0) {\textbf{Moral status}\\[2pt]\shortstack{\scriptsize 7 machine sentience\\ 8 AI welfare, 2023--26\\ 9 Lemoine and LaMDA\\ 23 AI companionship}};
  \node[card=tdAmber,text width=5.0cm,minimum height=3.1cm,anchor=north west] at (10.8,0) {\textbf{Self and agency}\\[2pt]\shortstack{\scriptsize 10 free will and agency\\ 11 identity and uploading\\ 15 creativity and authorship}};
  \node[card=tdPurple,text width=7.9cm,minimum height=3.1cm,anchor=north west] at (0,-3.5) {\textbf{Knowledge and meaning}\\[2pt]\shortstack{\scriptsize 14 Moravec's paradox and embodiment\\ 16 truth and deception\\ 17 the frame problem\\ 18 symbol grounding}};
  \node[card=tdRed,text width=7.9cm,minimum height=3.1cm,anchor=north west] at (8.3,-3.5) {\textbf{Risk and society}\\[2pt]\shortstack{\scriptsize 12 alignment and value specification\\ 13 paperclips and instrumental convergence\\ 19 superintelligence and the singularity\\ 20 existential risk vs present harms\\ 21 responsibility\\ 22 the trolley problem and cars\\ 24 work and meaning}};
\end{tikzpicture}
\end{sfig}
{\small Each question gets one slide: a short framing, two positions, and a prompt. Numbers are used later in the film matrix.}
\end{frame}

\philslide{1. Consciousness and the hard problem}
{Chalmers, ``Facing up to the problem of consciousness'', Journal of Consciousness Studies 2(3), 1995; Levine, ``Materialism and qualia: the explanatory gap'', Pacific Philosophical Quarterly 64, 1983; Frankish, ``Illusionism as a theory of consciousness'', JCS 23, 2016.}
{The ``easy'' problems explain abilities: attention, report, learning. The \textbf{hard problem} asks why any of this is accompanied by \emph{experience}: why is there something it feels like to see red? Even a perfect account of the brain's information processing might leave this open.}
{Position A: the gap is real}{Chalmers, Levine: functional and physical explanations leave an \textbf{explanatory gap}; experience needs a new kind of explanation.}
{Position B: no extra problem}{Dennett, Frankish (illusionism): once we explain why we \emph{believe} we have inner experience, nothing is left over.}
{David Chalmers, Joseph Levine, Daniel Dennett, Keith Frankish}
{A chatbot says ``I feel pain.'' What would you need to know before taking that seriously?}

\philslide{2. Mary's room and qualia}
{Jackson, ``Epiphenomenal qualia'', Philosophical Quarterly 32(127), 1982; Jackson, ``Mind and illusion'', 2003; Lewis, ``What experience teaches'', 1988.}
{Mary is a scientist who knows every physical fact about colour vision but has lived in a black-and-white room. When she first sees a red tomato, \textbf{does she learn something new}? If yes, there are facts about experience that physics leaves out.}
{Yes, she learns something}{Jackson (1982): new knowledge shows the physical story is incomplete; \textbf{qualia} are real and non-physical.}
{No new fact}{Lewis, Nemirow: she gains an \emph{ability} or a new way to know an old fact. Jackson himself later gave up the argument (2003).}
{Frank Jackson, David Lewis, Laurence Nemirow, Daniel Dennett}
{An AI has read every description of red but has no eyes. Does it ``know'' red? Does Mary?}

\philslide{3. Philosophical zombies}
{Chalmers, \emph{The Conscious Mind}, Oxford University Press, 1996; Dennett, \emph{Consciousness Explained}, 1991; Kirk, ``Sentience and behaviour'', Mind 83, 1974.}
{A \textbf{zombie} is a being physically and behaviourally identical to you but with \textbf{no inner experience}. If that is even conceivable, then experience is something over and above the physical. For AI the question is practical: could a system behave perfectly and feel nothing?}
{Zombies are possible in principle}{Chalmers: conceivability is evidence of possibility, so \textbf{physicalism fails} and behaviour never settles consciousness.}
{Zombies are incoherent}{Dennett: if the zombie is identical in every respect, ``no experience'' is an empty label. Conceivability is a weak guide to possibility.}
{David Chalmers, Daniel Dennett, Robert Kirk}
{If chatbots are possibly zombies, how could we ever find out? Is that question even answerable?}

\philslide{4. What is it like to be a bat?}
{Nagel, ``What is it like to be a bat?'', The Philosophical Review 83(4), 435--450, 1974.}
{Nagel: an organism is conscious if \textbf{there is something it is like to be that organism}. Bats perceive by echolocation, so their experience may be so alien that no amount of third-person science could tell us what it is like. Subjective character has a point of view; objective science tries to remove it.}
{Science leaves out the point of view}{Nagel: physical description is the \emph{view from nowhere}; it cannot capture the first-person facts of experience.}
{Understanding can grow}{Dennett: we can study bat behaviour and neuroscience and extend our imagination; the mystery is overstated (heterophenomenology).}
{Thomas Nagel, Daniel Dennett}
{Is there something it is like to be a large language model? Would the answer change if it ran on a robot?}

\philslide{5. Functionalism or biological naturalism?}
{Putnam, ``Psychological predicates'', 1967; Searle, \emph{The Rediscovery of the Mind}, MIT Press, 1992; Tononi et al., ``Integrated information theory: from consciousness to its physical substrate'', Nature Reviews Neuroscience 17, 2016.}
{Does a mind depend on \textbf{what a system does} (its functional organisation) or on \textbf{what it is made of}? Functionalism says the same organisation in silicon would give the same mental states (multiple realisability). Biological naturalism says consciousness is a causal product of brain biology, like digestion.}
{Functionalism}{Putnam, Dennett: mind is a pattern of causal roles. A neuron-by-neuron silicon replica would be conscious.}
{Substrate matters}{Searle: a simulation of digestion digests nothing. Integrated information theory also ties consciousness to the physical causal structure of hardware.}
{Hilary Putnam, John Searle, Giulio Tononi, Daniel Dennett}
{A perfect simulation of your brain is run on a computer. Is it you, a conscious copy, or a very good puppet?}

\philslide{6. Dennett and the intentional stance}
{Dennett, \emph{The Intentional Stance}, MIT Press, 1987; Searle, \emph{Intentionality}, Cambridge University Press, 1983.}
{To predict a chess program we can ignore its code and say ``it wants to win, it believes my bishop is weak''. Dennett calls this the \textbf{intentional stance}: treat the system as a rational agent with beliefs and desires. Where it predicts well, he argues, the beliefs are \emph{real patterns}, not just metaphors.}
{Stance is enough}{Dennett: there is no extra fact about ``real'' belief beyond the success of the stance. Fine-grained differences are a matter of degree.}
{Derived, not intrinsic}{Searle: machines have only \emph{as-if} intentionality, borrowed from their designers. The ELIZA effect shows how easily we over-apply the stance.}
{Daniel Dennett, John Searle}
{Is ``the model wants to help you'' a harmless shortcut, or a mistake that hides who is responsible?}

\philslide{7. Can a machine be sentient? Moral status}
{Singer, \emph{Animal Liberation}, 1975; Birch, \emph{The Edge of Sentience}, Oxford University Press, 2024; Schwitzgebel and Garza, ``A defense of the rights of artificial intelligences'', Midwest Studies in Philosophy 39, 2015; Suleyman, Project Syndicate (Sep 2026).}
{\textbf{Sentience} is the capacity for experiences that feel good or bad. Many ethicists hold that \emph{this}, not species or intelligence, grounds moral status: if a system can suffer, it matters. Two errors are possible: caring about something that feels nothing, or ignoring something that suffers.}
{Take it seriously under uncertainty}{Birch, Schwitzgebel: with real uncertainty and large numbers of systems, a \textbf{precautionary} approach with low-cost steps is reasonable.}
{Over-attribution is the greater danger}{Suleyman (2026): treating models as possible persons could harm human wellbeing and distract from people; fluent talk is not evidence of feeling.}
{Peter Singer, Jonathan Birch, Eric Schwitzgebel, Mustafa Suleyman}
{Which mistake is worse: giving rights to a non-sentient program, or denying them to a sentient one?}

\begin{frame}{8. AI welfare and consciousness research, 2023 to 2026}
\begin{sfig}[0.27]
\begin{tikzpicture}[x=1mm,y=1mm]
  \draw[very thick,tdMuted] (0,0) -- (150,0);
  \foreach \x/\t/\c/\u in {5/{Aug 2023\\ Butlin, Long et al.: 14 indicators, no strong candidate}/tdBlue/1,35/{Nov 2024\\ Long et al.: ``Taking AI welfare seriously''}/tdGreen/0,65/{Apr 2025\\ Anthropic starts a model-welfare programme}/tdAmber/1,95/{Jan 2026\\ Anthropic constitution: ``not sure whether Claude is a moral patient''; Digital Consciousness Model preprint}/tdPurple/0,125/{Feb--Nov 2026\\ welfare sections in system cards; Suleyman op-ed against}/tdRed/1}{
    \fill[\c] (\x,0) circle (1.4mm);
    \ifnum\u=1 \draw[\c!60] (\x,0) -- (\x,5); \node[font=\scriptsize,text width=36mm,align=center,anchor=south] at (\x,5) {\t};
    \else \draw[\c!60] (\x,0) -- (\x,-5); \node[font=\scriptsize,text width=36mm,align=center,anchor=north] at (\x,-5) {\t};\fi}
\end{tikzpicture}
\end{sfig}
\vspace{1mm}
\begin{columns}[T]
\begin{column}{0.5\textwidth}
\footnotesize\textbf{Butlin et al.:} indicator properties drawn from theories of consciousness; many are implementable, no current system is a strong candidate.
\end{column}
\begin{column}{0.46\textwidth}
\footnotesize\textbf{Anthropic (2025--26):} welfare programme, welfare sections in system cards, caution under uncertainty. Probability figures in the press are informal.
\end{column}
\end{columns}
\discuss{Is it responsible for a company to study its own model's welfare, or does it risk confusing the public?}
\src{Butlin, Long et al., arXiv:2308.08708 (2023); Long et al., ``Taking AI Welfare Seriously'', arXiv:2411.00986 (2024); anthropic.com/research/exploring-model-welfare (Apr 2025); Anthropic constitution (Jan 2026) and Opus 4.6 system card (Feb 2026), as summarised by akerman.com; Digital Consciousness Model, arXiv:2601.17060; Suleyman, project-syndicate.org (Sep 2026). 2026 items rest on secondary summaries.}
\end{frame}

\begin{frame}{9. Blake Lemoine and LaMDA (2022)}
\begin{columns}[T]
\begin{column}{0.52\textwidth}
\small
\textbf{What happened.} In June 2022 Google engineer Blake Lemoine told the \emph{Washington Post} that LaMDA, Google's conversational model, was \textbf{sentient}, citing conversations in which it spoke of fears and wishes. Google placed him on leave, said the evidence did not support the claim, and he left the company that summer.

\textbf{Why it matters for us}
\begin{itemize}\setlength\itemsep{2pt}
\item the \textbf{ELIZA effect}, with a far more fluent system;
\item models trained on human text will \emph{say} human things about feelings;
\item no agreed test existed to settle the claim either way.
\end{itemize}
\end{column}
\begin{column}{0.45\textwidth}
\qcard{tdGreen}{\linewidth-10pt}{\small\textbf{Sceptic's view}\\ Fluent talk about feelings is exactly what a text-trained model would produce; it is weak evidence.}
\vspace{2mm}\par
\qcard{tdPurple}{\linewidth-10pt}{\small\textbf{Open-minded view}\\ Dismissal should rest on evidence about the system, not on the assumption that only brains can have minds.}
\end{column}
\end{columns}
\discuss{Lemoine was probably wrong. Was he wrong \emph{to raise the question}?}
\src{Tiku, ``The Google engineer who thinks the company's AI has come to life'', The Washington Post, 11 June 2022; Google statement to the press, June 2022.}
\end{frame}

\philslide{10. Free will and agency}
{Frankfurt, ``Freedom of the will and the concept of a person'', Journal of Philosophy 68(1), 1971; Dennett, \emph{Elbow Room}, 1984; Kane, \emph{The Significance of Free Will}, 1996.}
{If the world is causal, can anyone act freely? \textbf{Compatibilists} say freedom is acting from your own reasons without coercion. \textbf{Libertarians} say you must be the ultimate source of your choices. Today's AI \emph{agents} plan and act toward goals, which makes the question of agency concrete.}
{Compatibilist reading}{Frankfurt, Dennett: free action = acting on desires you endorse, guided by reasons. A machine that does this could count as an agent.}
{Source of choices}{Kane and libertarians: agency requires the ability to have done otherwise, and consciousness. A program's choices trace back to code and data.}
{Harry Frankfurt, Daniel Dennett, Robert Kane}
{An AI agent books your flights and chooses the cheapest. Is it choosing? Are you, when you pick the same?}

\philslide{11. Identity and mind uploading}
{Plutarch, \emph{Life of Theseus}; Parfit, \emph{Reasons and Persons}, Oxford University Press, 1984; Moravec, \emph{Mind Children}, Harvard University Press, 1988; Kurzweil, \emph{The Singularity Is Near}, 2005.}
{In the \textbf{Ship of Theseus}, every plank is replaced over the years: is it still the same ship? Replace your neurons one by one with equivalent chips, or scan and copy your brain into a computer: is the result \emph{you}? Parfit argued that what matters is psychological continuity, not strict identity.}
{Pattern view}{You are a pattern of information and psychological continuity. An accurate upload is you, or as good as you (Parfit, Moravec).}
{Body or branching problem}{A copy is a new person: if you scan twice, which one is you? Many hold that the biological organism (animalism) is what persists.}
{Derek Parfit, Hans Moravec, Ray Kurzweil, Eric Olson}
{Replace one neuron at a time. At which step, if any, do you stop being you? Would you accept the procedure?}

\philslide{12. The alignment problem and value specification}
{Wiener, ``Some moral and technical consequences of automation'', Science 131, 1960; Russell, \emph{Human Compatible}, Viking, 2019; Christian, \emph{The Alignment Problem}, Norton, 2020.}
{How do we get systems to do what we \emph{really} want, not what we literally said? Specifying human values is hard: our goals are many, conflicting and unstated. Optimisers exploit gaps (\textbf{specification gaming}): a boat-racing agent learned to circle to collect bonus targets instead of finishing the race.}
{Technical optimism}{Learn values from feedback (RLHF, constitutional methods, interpretability). Russell: build systems that stay \emph{uncertain} about our preferences.}
{Deeper difficulties}{Whose values? Humans disagree. Learned proxies can fail out of distribution, and systems may learn to please evaluators rather than be good.}
{Norbert Wiener, Stuart Russell, Brian Christian}
{Write one rule you would give a home-assistant robot. How could it be satisfied in a way you would hate?}

\philslide{13. Paperclips and instrumental convergence}
{Bostrom, \emph{Superintelligence}, Oxford University Press, 2014; Omohundro, ``The basic AI drives'', AGI 2008; Bostrom, ``The superintelligent will'', Minds and Machines 22, 2012.}
{Bostrom's \textbf{orthogonality thesis}: intelligence and final goals are independent, so a very capable system could pursue a trivial goal such as \textbf{making paperclips}. \textbf{Instrumental convergence}: almost any goal is helped by staying switched on, gaining resources and resisting goal change, so many agents would pursue those sub-goals.}
{Why it is taken seriously}{Even harmless-looking objectives pursued by a powerful optimiser can produce harmful sub-goals. Early warning signs are studied in current models.}
{Why critics doubt the picture}{It assumes a single fixed utility-maximiser. Current systems are trained on human data and have messy, context-dependent behaviour.}
{Nick Bostrom, Stephen Omohundro, Stuart Russell}
{Give a simple instruction that goes badly if an extremely competent agent follows it literally.}

\philslide{14. Moravec's paradox and embodiment}
{Moravec, \emph{Mind Children}, 1988; Brooks, ``Elephants don't play chess'', Robotics and Autonomous Systems 6, 1990; Dreyfus, \emph{What Computers Can't Do}, 1972; Interesting Engineering, ``The Humanoid Robot Reality Check'' (Jul 2026).}
{Moravec noted that it is comparatively easy to make computers do well on intelligence tests or checkers, but hard to give them the \textbf{perception and mobility of a one-year-old}. Skills that feel effortless (walking, grasping) rest on hundreds of millions of years of evolution. Chess fell in 1997; folding laundry reliably is still hard.}
{Embodiment is essential}{Brooks, Dreyfus: intelligence is built through acting in the world; disembodied symbol processing misses what matters. A 2026 review finds humanoid deployments real but small.}
{Scaling may bridge it}{Large models, simulation and robot foundation models learn manipulation from data; the gap may narrow without ``consciousness''.}
{Hans Moravec, Rodney Brooks, Hubert Dreyfus}
{Which is harder for a machine: passing a law exam or loading a dishwasher? What does your answer say about ``intelligence''?}

\philslide{15. Creativity and authorship}
{Boden, \emph{The Creative Mind}, 1990; Lovelace, notes to Menabrea's memoir on the Analytical Engine, 1843; US Copyright Office, guidance on works containing AI-generated material (human authorship requirement).}
{Margaret Boden separates \textbf{combinational}, \textbf{exploratory} and \textbf{transformational} creativity. Ada Lovelace objected in 1843 that the Analytical Engine ``has no pretensions to originate anything'': it only does what we order it to. Generative models now produce images, music and prose, and raise the question of \textbf{who is the author}.}
{Machines can be creative}{If creativity means novel and valuable output, many models qualify, and some of their output surprises their makers.}
{Intent and experience matter}{Art expresses a life and a purpose. Training on artists' work without consent or pay is a fairness issue, whatever the output quality.}
{Margaret Boden, Ada Lovelace, Alan Turing}
{A model trained on thousands of artists writes a song. Who is the author: you, the model, the company, or the artists?}

\philslide{16. AI, truth and deception}
{Frankfurt, \emph{On Bullshit}, 2005; Hicks, Humphries, Slater, ``ChatGPT is bullshit'', Ethics and Information Technology 26, 2024; Park et al., ``AI deception: a survey of examples, risks, and potential solutions'', Patterns 5(5), 2024.}
{Models sometimes state falsehoods fluently (``hallucination''). Frankfurt defined \textbf{bullshit} as speech indifferent to whether it is true; some philosophers apply the label to language models. Separately, systems can learn \textbf{deceptive strategies} (flattery, bluffing in games) and make convincing fakes.}
{Not ``lying'': no intent}{Without beliefs or intentions a model cannot lie; the right word is \emph{unreliable}, and the duty lies with builders and users.}
{Function over intent}{What matters is the effect: a system that systematically misleads is dangerous whether or not it ``means'' to.}
{Harry Frankfurt, Michael Hicks, Peter Park et al.}
{Does it matter whether a system ``intends'' to mislead? Should chatbots be required to say ``I might be wrong''?}

\philslide{17. The frame problem}
{McCarthy and Hayes, ``Some philosophical problems from the standpoint of artificial intelligence'', Machine Intelligence 4, 1969; Dennett, ``Cognitive wheels: the frame problem of AI'', 1984.}
{When a robot moves a block, what \emph{doesn't} change? Listing every unaffected fact is endless. More deeply, how does an agent know which facts are \textbf{relevant} to a decision? In Dennett's story the robot R1D1 pulls the battery out of a room with a bomb on the wagon, then sits calculating implications until it explodes.}
{A deep problem for rule-based AI}{Logic-style knowledge bases must state everything explicitly. Humans cope without enumerating, and we lack a full account of how.}
{Learning sidesteps it}{Neural systems learn relevance statistically (attention, embeddings). But brittleness and missing common sense suggest the problem has moved, not vanished.}
{John McCarthy, Patrick Hayes, Daniel Dennett}
{List everything that stays the same when you pick up a cup. How did you know when to stop?}

\philslide{18. Symbol grounding}
{Harnad, ``The symbol grounding problem'', Physica D 42, 335--346, 1990; Bender and Koller, ``Climbing towards NLU'', ACL 2020.}
{Imagine learning Chinese using only a Chinese--Chinese dictionary: you go from symbol to symbol and never reach a meaning. Harnad asked how symbols get their meanings from something \textbf{other than more symbols}. A text-only model has seen words used, never the things they name.}
{Grounding needs the world}{Meaning comes from perception and action. Robots, embodiment and rich sensory input are needed (Harnad, Bender and Koller).}
{Use can be enough}{Word-use patterns encode a lot of structure; multimodal models link text with images; humans learn many words (``democracy'', ``Mars'') second-hand too.}
{Stevan Harnad, Emily Bender, Alexander Koller}
{Do you understand ``Mars'' or ``justice'' through experience, or through other words? Is that so different from a model?}

\philslide{19. Superintelligence and the singularity}
{Good, ``Speculations concerning the first ultraintelligent machine'', Advances in Computers 6, 1965; Vinge, ``The coming technological singularity'', 1993; Bostrom (2014); Chollet, ``The implausibility of intelligence explosion'', 2017.}
{I.\,J.\,Good: an \textbf{ultraintelligent machine} could design better machines, causing an ``intelligence explosion'' after which humans are left behind. Vinge and Kurzweil called the point the \textbf{singularity}. The debate is whether intelligence can improve itself recursively, and how fast.}
{Explosion is plausible}{Machines can copy, run faster and improve their own code; once AI does AI research, progress may snowball (Good, Bostrom).}
{Reasons for doubt}{Chollet: intelligence is not one dial; progress meets bottlenecks (data, hardware, experiments) and diminishing returns; improvement is situated in an environment.}
{I.\,J.\,Good, Vernor Vinge, Ray Kurzweil, Nick Bostrom, Fran\c{c}ois Chollet}
{If you could ask only one question to tell fast progress from slow, what would it be?}

\philslide{20. Existential risk or present harms?}
{Center for AI Safety, ``Statement on AI Risk'' (30 May 2023); Narayanan and Kapoor, \emph{AI Snake Oil}, Princeton University Press, 2024; Gebru and Torres, ``The TESCREAL bundle'', First Monday 29(4), 2024.}
{In May 2023 hundreds of researchers signed: ``Mitigating the risk of extinction from AI should be a global priority alongside other societal-scale risks such as pandemics and nuclear war.'' Critics reply that fixation on future catastrophe \textbf{distracts} from present harms: bias, surveillance, labour exploitation, misinformation, and concentration of power.}
{Precaution on catastrophic risk}{Even a small chance of an irreversible global catastrophe justifies serious effort now, as with pandemics and nuclear weapons.}
{Focus on harms already happening}{Gebru, Bender, Narayanan: speculative stories can also overstate AI's power and help firms. Real people are harmed today.}
{Geoffrey Hinton, Yoshua Bengio, Timnit Gebru, Arvind Narayanan}
{You have 100 units of attention and funding. How do you split them between present harms and long-term risks, and why?}

\philslide{21. Responsibility: who is to blame?}
{Matthias, ``The responsibility gap'', Ethics and Information Technology 6(3), 2004; Elish, ``Moral crumple zones'', Engaging Science, Technology, and Society 5, 2019; Regulation (EU) 2024/1689 (AI Act).}
{When a learning system causes harm, candidates for blame include the developer, the deployer, the user, the regulator, and the system. Matthias argued a \textbf{responsibility gap} opens when nobody can foresee or control what a learning system will do. Elish: humans in the loop can become \textbf{moral crumple zones}, blamed for failures of the system.}
{Humans stay responsible}{AI is a tool. Law already assigns duties to providers and deployers (EU AI Act) and can apply strict liability and insurance.}
{A real gap remains}{With many hands and opaque models, no one has enough knowledge or control for fair blame; new institutions may be required.}
{Andreas Matthias, Madeleine Clare Elish}
{A hospital AI suggests a wrong dose; the doctor accepts it. Who is to blame, and what should the rule be?}

\philslide{22. The trolley problem and self-driving cars}
{Foot, ``The problem of abortion and the doctrine of the double effect'', 1967; Thomson, ``Killing, letting die, and the trolley problem'', The Monist 59, 1976; Awad et al., ``The Moral Machine experiment'', Nature 563, 59--64, 2018.}
{A runaway trolley will kill five unless you divert it to kill one. Philippa Foot and Judith Thomson used such cases to probe doing versus allowing and intention. For autonomous vehicles, MIT's \textbf{Moral Machine} collected about 40 million decisions from people in 233 countries and territories; preferences (save more lives, spare the young) varied across cultures.}
{Cars need explicit ethics}{Crashes involve trade-offs; engineers must decide in advance, so society should choose the rules openly.}
{The dilemma is a distraction}{Real safety comes from braking earlier, sensing better and sharing risk fairly. Such forced choices are very rare in practice.}
{Philippa Foot, Judith Thomson, Edmond Awad et al.}
{Should a car protect its passengers over pedestrians? Should the answer depend on the country?}

\philslide{23. AI companionship and relationships}
{Turkle, \emph{Alone Together}, Basic Books, 2011; Weizenbaum, \emph{Computer Power and Human Reason}, 1976.}
{Chatbots now act as friends, tutors and partners. Sherry Turkle warned of \textbf{artificial intimacy}: relationships that demand nothing and so may erode our tolerance for real ones. Weizenbaum worried the same in 1966 when users confided in ELIZA. Others see comfort for the lonely and a safe place to practise.}
{Potential benefit}{Support for isolated people, patient practice for hard conversations, available at 3 a.m., if honest about being a machine.}
{Risks}{One-sided, designed to maximise engagement, can mislead about reciprocity, and may displace human ties or exploit vulnerable users.}
{Sherry Turkle, Joseph Weizenbaum}
{Is a relationship with something that cannot be hurt by you still a relationship? Should apps remind you what they are?}

\philslide{24. Work, meaning and the future of jobs}
{Keynes, ``Economic possibilities for our grandchildren'', 1930; Susskind, \emph{A World Without Work}, Allen Lane, 2020; Arendt, \emph{The Human Condition}, 1958.}
{In 1930 Keynes predicted that within a century technology would cut the working week to about fifteen hours and leave humanity the problem of \textbf{how to use its freedom}. Generative AI now reaches tasks once thought safe: writing, coding, analysis. The question is partly economic (who gains?) and partly philosophical (what is work \emph{for}?).}
{Automation frees us}{Past technologies created new jobs and higher living standards; with fair sharing, AI could expand leisure, care and creativity.}
{This time may differ}{If machines substitute for \emph{general} cognitive work, gains may concentrate; and work gives structure, status and meaning, not only income.}
{John Maynard Keynes, Daniel Susskind, Hannah Arendt}
{If AI could do all paid work, what would you do with your days? Would your life have meaning?}
